\chapter{Fourier, filtre și convoluție: instrumentele de bază}
\label{ch:fundamente}

Restul documentului se sprijină pe trei idei simple, în ordine: orice semnal
e o sumă de frecvențe, un filtru alege care dintre ele rămân, iar convoluția
e mecanismul care face asta. Wavelets, câteva capitole mai încolo, e ce se
întâmplă când combini toate trei cu o singură idee nouă.

\section{Transformata Fourier}

Shazam recunoaște o piesă din zece secunde de audio pentru că amprenta ei de
frecvențe e unică~\cite{wang2003}, adică exact amprenta pe care o calculează
transformata Fourier (revenim pe larg la Shazam în
Capitolul~\ref{ch:shazam}). Rezultatul central al transformatei e că orice
semnal se scrie ca o sumă, sau pentru semnale continue o integrală, de
sinusoide~\cite{lyons2011}:
\[
F(\omega) = \int_{-\infty}^{\infty} f(t)\,e^{-i\omega t}\,dt, \qquad
f(t) = \frac{1}{2\pi}\int_{-\infty}^{\infty} F(\omega)\,e^{i\omega t}\,d\omega,
\]
unde pulsația $\omega = 2\pi f$ e aceeași mărime ca frecvența $f$, doar în alte
unități (rad/s față de Hz). Formula lui Euler, $e^{i\theta} = \cos\theta +
i\sin\theta$, e ce leagă exponențiala complexă de sinusoide. Teorema lui
Parseval,
\[
\int_{-\infty}^{\infty} |f(t)|^2\,dt = \frac{1}{2\pi}\int_{-\infty}^{\infty} |F(\omega)|^2\,d\omega,
\]
spune că energia semnalului e aceeași, calculată în timp sau în frecvență. E o
proprietate la care revenim când arătăm că nici transformata wavelet
discretă nu pierde energie.

Limita transformatei contează la fel de mult ca rezultatul: Fourier spune CE
frecvențe conține un semnal, dar nu CÂND apar ele. Un mashup din două piese
are, pe porțiuni lungi, un spectru foarte asemănător cu suprapunerea lor,
tocmai de-aceea Shazam rulează transformata pe ferestre scurte și succesive,
nu pe piesa întreagă (transformata Fourier pe termen scurt, STFT).
Fereastra scurtă fixează unde se pierde: cu cât fereastra e mai scurtă, cu
atât localizarea în timp e mai bună și rezoluția în frecvență mai slabă,
exact compromisul pe care wavelets îl rezolvă altfel, cu ferestre de
dimensiune variabilă (Capitolul~\ref{ch:wavelets}).

Demo-ul din tur (Figura~\ref{fig:fourier-demo}) calculează efectiv o
transformată Fourier rapidă (\txt{np.fft.fft}) pentru opt funcții
predefinite: sinusoidă simplă, două și trei frecvențe, chirp, puls
Gaussian, undă pătrată, suma Fourier a undei pătrate și o expresie
personalizată, evaluată printr-o listă restrânsă de funcții permise
(\txt{sin}, \txt{cos}, \txt{exp}, \txt{pi}, \txt{abs}, \txt{sqrt}). Afișează
simultan domeniul timp și spectrul de magnitudine.

Presetul cu suma Fourier a undei pătrate arată o limită care revine mai
târziu. Cu fiecare armonică adăugată suma se apropie de unda pătrată, dar
lângă discontinuitate rămâne o supraoscilație de aproape nouă la sută din
saltul semnalului. Ea se îngustează pe măsură ce adaugi termeni și nu scade
niciodată în înălțime. Fenomenul poartă numele lui Gibbs și e ce se întâmplă
ori de câte ori tai brusc un spectru. Pentru presetul
din figură, $f(t) = \sin(2\pi \cdot 5 t) + \sin(2\pi \cdot 12 t)$, cele două
vârfuri spectrale cad exact la 5 și 12~Hz.

\begin{figure}[htbp]
\centering
\includegraphics[width=\textwidth]{05-fourier-demo.png}
\caption{Domeniul timp și spectrul de magnitudine calculate live pentru
  presetul „Două frecvențe”: vârfurile spectrale cad exact la 5 și 12~Hz,
  frecvențele din formula afișată dedesubt.}
\label{fig:fourier-demo}
\end{figure}

\section{Filtre digitale}
\label{sec:filtre}

Un filtru păstrează o parte din spectru și atenuează restul, teoria
standard de proiectare a filtrelor digitale~\cite{oppenheimschafer2010}.
Trei forme apar în tur, toate normalizate să treacă prin $-3$~dB (adică
$|H|=1/\sqrt2$, jumătate din putere) exact la frecvența de tăiere $f_c$, ca
să fie comparabile la aceeași tăietură:
\begin{gather*}
H^{\text{ideal}}_{LP}(f) = \begin{cases} 1 & |f|\le f_c \\ 0 & |f| > f_c \end{cases} \\[4pt]
\left|H^{\text{Butterworth}}_{LP}(f)\right|^2 = \frac{1}{1+\left(f/f_c\right)^{2n}} \\[4pt]
H^{\text{Gaussian}}_{LP}(f) = e^{-\frac{\ln 2}{2}\left(f/f_c\right)^2}
\end{gather*}
Filtrul ideal are un prag dur chiar la $f_c$, imposibil de realizat în
practică: răspunsul lui în timp e un \emph{sinc} infinit, deci necauzal, iar
trunchierea lui readuce oscilația Gibbs întâlnită deja la Fourier. Butterworth
și Gaussian sunt calibrate să dea exact $|H|^2=\tfrac12$ la $f=f_c$, oricare
ar fi ordinul $n$ sau lățimea Gaussianei. Ordinul controlează doar cât de
abrupt coboară restul curbei, de la $n=1$ (lent) la $n=8$ (aproape ideal).
Varianta trece-sus e complementul în putere al variantei trece-jos,
\[
H_{HP}(f) = \sqrt{1-\left|H_{LP}(f)\right|^2}, \qquad
\text{deci} \quad |H_{LP}|^2+|H_{HP}|^2=1
\]
la orice frecvență: cele două filtre împart mereu toată puterea, indiferent
de formă sau ordin. Varianta trece-bandă înlănțuie un trece-sus la tăietura
joasă cu un trece-jos la tăietura înaltă, $H_{BP}(f) =
\sqrt{1-|H_{LP}(f;f_{lo})|^2}\cdot H_{LP}(f;f_{hi})$.

\begin{figure}[htbp]
\centering
\includegraphics[width=\textwidth]{07-filters-theory.png}
\caption{Cele trei forme ale filtrului trece-jos, toate normalizate la
  $-3$~dB pe frecvența de tăiere $f_c$.}
\label{fig:filters-theory}
\end{figure}

În back-end, o singură funcție, \txt{filter_magnitude}, produce atât curba
de răspuns desenată pe ecran (Figura~\ref{fig:filters-demo}), cât și
multiplicatorul aplicat efectiv semnalului: cererea
\txt{/api/filters/apply-signal} transformă semnalul cu FFT, îl înmulțește cu
acel multiplicator și aplică FFT inversă. Multiplicatorul e real și par în
frecvență, așa că semnalul filtrat rămâne real și filtrul e zero-fază. Iar
faptul că e o singură funcție înseamnă că demo-ul nu poate afișa un filtru
și calcula altul, o clasă întreagă de bug posibil eliminată prin construcție.

\begin{figure}[htbp]
\centering
\includegraphics[width=\textwidth]{08-filters-demo.png}
\caption{Filtru trece-jos Butterworth de ordinul 4: răspunsul în frecvență
  (stânga sus), spectrul semnalului înainte și după filtrare (dreapta sus) și
  semnalul filtrat în timp (jos), toate recalculate la orice schimbare de
  parametru.}
\label{fig:filters-demo}
\end{figure}

Băncile de filtre, cu un trece-jos $h$ pentru coeficienți de aproximare și un
trece-sus $g$ pentru coeficienți de detaliu, aplicate recursiv, dau analiza
multi-rezoluție~\cite{strangnguyen1996}, punctul de plecare al
Capitolului~\ref{ch:wavelets}. Condiția care garantează reconstrucție
perfectă e relația de filtre-oglindă-conjugate,
$g[k] = (-1)^k h[N-1-k]$~\cite{estebangaland1977,smithbarnwell1984}, iar
motivul pentru care decimarea cu doi după fiecare filtrare nu pierde
informație vine direct din teorema eșantionării: o bandă de $B$~Hz cere
doar $2B$ eșantioane pe secundă~\cite{shannon1949}, exact atât cât rămâne
după ce fiecare filtru înjumătățește banda.

\section{Convoluția}

Definiția 1D, $(f*g)[n] = \sum_k f[k]\,g[n-k]$, stă la baza filtrelor de mai
sus și a rețelelor convoluționale deopotrivă. Teorema convoluției,
$f*g \leftrightarrow F\cdot G$~\cite{lyons2011}, leagă cele două: o
convoluție în timp e o simplă înmulțire în frecvență, motivul pentru care
FFT reduce complexitatea unei convoluții lungi de la $O(n^2)$ la
$O(n\log n)$: transformi, înmulțești punct cu punct, transformi înapoi.

Extensia la imagini,
\[
(I*K)[x,y] = \sum_{i,j} I[x-i,\,y-j]\cdot K[i,j],
\]
rotește kernelul cu $180^\circ$ (semnele minus). Bibliotecile de procesare
de imagini, și demo-urile din tur, calculează de fapt \emph{corelație}, cu
$I[x+i,y+j]$, fără rotația kernelului: pentru kernele simetrice (blur,
Gaussian, Laplacian) rezultatul e identic, iar diferența apare doar la
kernele asimetrice, la Sobel de exemplu, unde rotația ar inversa semnul
gradientului (Figura~\ref{fig:conv-2d}).

\begin{figure}[htbp]
\centering
\includegraphics[width=\textwidth]{12-conv-2d.png}
\caption{Convoluția 2D cu semnul corect și nota despre corelație versus
  convoluție, alături de trei kernele reprezentative: blur Gaussian, sharpen
  și Sobel X (asimetric).}
\label{fig:conv-2d}
\end{figure}

\section{Fourier în două dimensiuni}
\label{sec:fourier-2d}

Înainte de a trece la kernele, turul face un pas lateral care merită explicat,
pentru că e locul unde tot ce s-a spus despre semnale 1D se mută pe imagini.
Transformata Fourier discretă se aplică la fel de bine unei matrice de pixeli,
pe rânduri și apoi pe coloane, iar rezultatul e tot o matrice, de aceeași
dimensiune, pe care o vedem ca imagine: spectrul de magnitudine
(Figura~\ref{fig:sound-to-image}). Backend-ul o calculează la cerere, prin
\txt{/api/fourier/image}, pe imaginea aleasă din listă.

Se citește altfel decât un spectru 1D. Frecvențele joase stau în centru, cele
înalte spre margini, iar distanța față de centru înseamnă cât de repede se
schimbă intensitatea, nu unde se întâmplă schimbarea. O zonă întinsă și
uniformă, cerul sau un perete, produce energie concentrată în centru. Muchiile și
textura fină împing energie spre margini, iar orientarea se păstrează: o
textură cu dungi verticale lasă o urmă orizontală în spectru, perpendiculară
pe dungi. Scara e logaritmică, altfel componenta continuă din centru ar fi
singurul lucru vizibil.

Consecința practică leagă capitolul: un blur de telefon e un filtru trece-jos
în două dimensiuni, iar o accentuare de contur e un trece-sus, exact aceleași
două operații din Secțiunea~\ref{sec:filtre}, doar că acum tăietura se face pe
un disc în jurul centrului, nu pe un interval de pe o axă. Kernelele din
secțiunea următoare sunt fix aceste filtre, scrise în domeniul spațial.

\begin{figure}[htbp]
\centering
\includegraphics[width=\textwidth]{13-sound-to-image.png}
\caption{Aceeași imagine în cele două domenii: originalul în nivele de gri
  (stânga) și spectrul lui de magnitudine 2D pe scară logaritmică (dreapta),
  cu frecvențele joase în centru. Blana deasă a babuinului umple spectrul spre
  margini, semnătura unei texturi fine.}
\label{fig:sound-to-image}
\end{figure}

\section{Kernele 2D, pe familii}
\label{sec:kernele}

Un demo educațional descompune aplicarea unui kernel pas cu pas, cu fereastra
care alunecă, produsul element-cu-element și suma, înainte de a trece la
kernele aplicate pe imagini reale (Figura~\ref{fig:kernels-demo}): blur
(cutie sau Gaussian), ascuțire (normală sau puternică), detectare de muchii
(Laplacian, Laplacian diagonal, Sobel, Prewitt, contur) și emboss, fiecare
cu raza (3, 4 sau 5) și puterea efectului reglabile independent.

\begin{figure}[htbp]
\centering
\includegraphics[width=\textwidth]{16-kernels-demo.png}
\caption{Blur Gaussian $3\times 3$ aplicat pe imaginea Peppers, cu matricea
  de kernel afișată exact așa cum e trimisă către back-end.}
\label{fig:kernels-demo}
\end{figure}

Redimensionarea unui kernel la o rază mai mare trebuie să păstreze \emph{ce
face} kernelul, nu doar dimensiunea lui, iar o singură regulă nu ajunge
pentru toate familiile. Identitatea rămâne o deltă discretă la orice
dimensiune. Kernele de blur adevărat (toate valorile nenegative, suma 1)
devin cutie sau Gaussian, după cum sursa era uniformă sau nu. Kernele de
muchie (suma 0) păstrează tiparul $3\times 3$ pe inelul exterior și
rebalansează centrul ca suma să rămână 0 la orice dimensiune. Kernele
direcționale (emboss) sunt interpolate biliniar și rescalate ca să
păstreze suma originală, care le fixează luminozitatea rezultatului. Kernele
de ascuțire simetrică, cu suma 1 dar și valori negative, deci nu blur, se
reconstruiesc ca \emph{unsharp masking} la noua rază,
\[
K_{\text{sharpen}}(r) = (1+\alpha)\,\delta_r - \alpha\, G_r,
\]
unde $\delta_r$ e delta discretă, $G_r$ Gaussiana normalizată la raza $r$,
iar $\alpha$ vine din valoarea centrală a kernelului original minus 1, ca un
sharpen la $3\times 3$ și unul la $5\times 5$ să păstreze aceeași putere a
efectului.

Distincția pe familii a intrat în cod după ce o regulă mai simplă,
„orice kernel cu suma 1 e un blur”, s-a dovedit greșită: și ascuțirea, și
emboss-ul, și identitatea au suma 1, așa că un „Sharpen” cerut la
$5\times 5$ întorcea de fapt o gaussiană, cu un câștig măsurat de zero la
frecvența Nyquist. Cu funcția rescrisă pe familii, pe o fotografie reală,
varianța laplacianului imaginii, o măsură standard de claritate, urcă de
la 2307 pe original la 40150, 33893 și respectiv 32114 pentru ascuțire la
$3\times 3$, $5\times 5$ și $7\times 7$.
